\chapter{实验设计}
\label{chap:exp}

本章给出实验方案。需要强调的是，本文的重点是\textbf{系统设计}而非算法上的最优——采用标准 GRPO 只是为了\textbf{证明整套自进化系统的可行性}，而不追求某个强化学习算法细节上的 SOTA。因此本文\textbf{不设算法消融实验}，实验围绕"系统能否完整运作、多智能体协作能否抑制奖励作弊、奖励机制能否增强模型表现"三点展开。

\section{任务与环境}
\label{sec:exp-setup}

\textbf{任务}取自 ClawEval 的纯文本子集，共 195 个工具使用任务，覆盖工作流编排、系统操作、问答检索、财务、办公文档、代码等多个能力域。每个任务在代码沙箱中执行，以沙箱内的真实状态为完成与否的判据。\textbf{基座}为 Qwen3.5-9B，采用完全异步的策略分离部署、BF16 全栈。\textbf{策略优化}采用标准 GRPO。

\section{以少量种子查询驱动的训练}
\label{sec:exp-train}

本文进行一次 \textbf{20 步}的强化学习训练，用以验证整套系统能被真正"驱动起来"。其关键设计是：在固定的每步总轨迹预算下，\textbf{仅使用 $1/(1+\text{追问次数})$ 的种子查询数}，靠多轮追问把每个种子扩展为一段多轮对话，补足总轨迹量。

设每步组大小为 $G$、追问次数为 $K$，则单轮方案需要 $Q_0$ 个种子查询得到 $Q_0\cdot G$ 条轨迹；本文多轮方案只取 $Q_0/(1+K)$ 个种子，每个种子经 $1+K$ 轮交互，同样得到 $\frac{Q_0}{1+K}\cdot(1+K)\cdot G = Q_0\cdot G$ 条轨迹。\textbf{总量不变，但采样从单轮变为多轮}——用更少的种子、更深的交互，驱动同等规模的强化学习更新。这既贴近真实的多轮使用场景，也直接体现了多智能体采样架构（尤其是提问者）的价值。

\textbf{说明}：追问机制能否启用，取决于多轮采样链路在系统级验证中是否跑通（第\ref{sec:exp-e2e}节）。若多轮链路已稳定，则按上式以少量种子驱动训练；若受工程约束暂未跑通，则退化为单轮采样的等价总量训练，系统的其余部分（差分驱动奖励、GRPO 更新）不受影响。

\section{系统端到端验证}
\label{sec:exp-e2e}

验证分两级：

\begin{enumerate}
  \item \textbf{单元级}：采样、观察者取证、奖励聚合、轨迹装配等模块以 CPU 单元测试覆盖，确保各环节在无 GPU 环境下逻辑正确。
  \item \textbf{系统级}：在真实多卡环境上跑通完整训练步——采样产出有效轨迹、观察者产出状态差分、裁判正常打分、GRPO 完成策略更新。以\textbf{能否稳定完成若干训练步}、每步\textbf{有效轨迹产出率}、以及\textbf{奖励/优势离散度}是否正常，作为系统完整运作的判据。
\end{enumerate}

\section{多智能体协作对奖励作弊的抑制}
\label{sec:exp-hacking}

本文的核心论点之一是：\textbf{多智能体协作}——尤其是观察者以环境真实状态差分取证——能抑制奖励作弊。为此，统计执行者\textbf{自我声明}与观察者\textbf{状态差分}不一致的样本比例，考察差分驱动的裁判能否在这些样本上给出与真实状态一致（而非被自述误导）的奖励。这是"多 Agent 协作防 reward hacking"的直接证据。

\section{奖励机制对模型表现的增强}
\label{sec:exp-gain}

在系统完整运作的前提下，本文比较训练前后策略在任务上的表现（以 Pass$^N$ 稳定成功率刻画），验证差分驱动的奖励能否为策略提供有效的学习信号、从而\textbf{增强模型表现}。这里的目标不是刷高绝对分数，而是证明\textbf{整套自进化闭环确实在朝着更好的策略推进}。

\section{实验结果}

\todo{填入：系统端到端稳定运转的证据（有效轨迹产出率、reward/优势/组内 std 训练曲线）；20 步训练前后模型表现对比；执行者自述与观察者状态差分不一致样本上的抗作弊统计。当前 16 卡正式训练进行中，结果就绪后补入。}
